% Intended LaTeX compiler: pdflatex
\documentclass[11pt]{article}
\usepackage[utf8]{inputenc}
\usepackage[T1]{fontenc}
\usepackage{graphicx}
\usepackage{longtable}
\usepackage{wrapfig}
\usepackage{rotating}
\usepackage[normalem]{ulem}
\usepackage{amsmath}
\usepackage{amssymb}
\usepackage{capt-of}
\usepackage{hyperref}
\usepackage[newfloat]{minted}
\author{root}
\date{\today}
\title{}
\hypersetup{
 pdfauthor={root},
 pdftitle={},
 pdfkeywords={},
 pdfsubject={},
 pdfcreator={Emacs 29.3 (Org mode 9.6.15)}, 
 pdflang={English}}
\begin{document}

\tableofcontents

\section{H}
\label{sec:orgc8fa02f}
\begin{minted}[linenos,firstnumber=1,bgcolor=gray,frame=lines]{python}
def f():
    return 1
\end{minted}

\begin{minted}[linenos,firstnumber=10,bgcolor=gray,frame=lines]{python}
a = 1
b = 2
\end{minted}

\begin{minted}[linenos,firstnumber=16,bgcolor=gray,frame=lines]{python}
c = 3
\end{minted}

\begin{minted}[linenos,firstnumber=1,bgcolor=gray,frame=lines]{common-lisp}
(save-excursion
  (goto-char (point-min))
\end{minted}
See 1 and 2 and the first.

\begin{minted}[linenos,firstnumber=1,bgcolor=gray,frame=lines]{common-lisp}
(foo)      (a)
(bar)      (b)
\end{minted}
Ref a

\begin{minted}[linenos,firstnumber=1,bgcolor=gray,frame=lines]{common-lisp}
(x)      (kx)
\end{minted}
1

\begin{minted}[bgcolor=gray,frame=lines]{common-lisp}
(y)      (ly)
\end{minted}
ly

\begin{verbatim}
1  one
2  two
\end{verbatim}

\begin{minted}[linenos,firstnumber=1,bgcolor=gray,frame=lines]{python}
l1
l2
l3
l4
l5
l6
l7
l8
l9
l10
\end{minted}

\begin{minted}[linenos,firstnumber=3,bgcolor=gray,frame=lines]{common-lisp}
(indented)      (ind)
  (more)
\end{minted}
Look at line ind here.
\end{document}
